
\chapter*{Kurzfassung}
\addcontentsline{toc}{chapter}{Kurzfassung}

Die Stichprobenanalyse im Internal Audit erfordert die manuelle Prüfung zahlreicher Rechnungen und Freigaben und ist entsprechend zeitaufwendig. \acp{LLM} können diese dokumentenbasierte Prüfung unterstützen. Ihre Neigung zu Halluzinationen und nicht deterministischen Ausgaben wirft im Audit jedoch Fragen nach Zuverlässigkeit und Reproduzierbarkeit auf. Die vorliegende Arbeit untersucht daher, wie sich verschiedene LLMs hinsichtlich ihrer Eignung zur Unterstützung von Stichprobenanalysen im Auditbereich der Freudenberg Group (im Folgenden Freudenberg) unterscheiden.

Hierzu wurden die acht in Microsoft Copilot Studio verfügbaren Modelle von OpenAI (GPT-4.1, GPT-5, GPT-5.3, GPT-5.5) und Anthropic (Claude Sonnet 4.6, Opus 4.6, Opus 4.7, Opus 4.8) jeweils als Agent mit identischen, an den Prüfschritten des Corporate Audit ausgerichteten Instruktionen konfiguriert. Als Testdatensatz dienen 100 synthetisch erzeugte Rechnungen mit zugehörigen Freigabenachweisen. In diese wurden 24 Audit Exceptions aus sechs Fehlerkategorien eingebaut. Bewertet wurden die Modelle über eine Nutzwertanalyse mit den Kriterien Recall (40 \%), Precision (30 \%) und Konsistenz (30 \%). Die Konsistenz wurde anhand von fünf Durchläufen über 30 Rechnungen gemessen und die Robustheit der Gewichtung über eine Sensitivitätsanalyse geprüft.

Den höchsten Nutzwert erreicht Claude Opus 4.7 mit 92,9 von 100 Punkten (Recall 91,7\ \%, Precision 95,7 \%), gefolgt von Opus 4.8 mit 89,6 Punkten. Nur diese beiden Modelle erzielen bei Recall, Precision und Konsistenz durchgehend Werte über 80 \%. Die Modelle von Anthropic belegen die Plätze eins bis vier und liegen bei Recall und Konsistenz vor allen Modellen von OpenAI, die ein unausgewogenes Fehlerbild zeigen. GPT-4.1 übersieht fast die Hälfte der Exceptions, während bei GPT-5, GPT-5.3 und GPT-5.5 mehr als die Hälfte der Meldungen unberechtigt ist. Die Rangfolge bleibt in allen Szenarien der Sensitivitätsanalyse stabil. Neuere Modellversionen schneiden dabei nicht automatisch besser ab als ältere.

Für Freudenberg wird der Einsatz von Opus 4.7 oder Opus 4.8 als Agent zur Vorprüfung empfohlen. Da auch das beste Modell nicht alle Exceptions erkennt, dienen die Ergebnisse als Entscheidungshilfe, und die abschließende Bewertung verbleibt beim Auditor.
